% !TEX root = ../../thesis.tex
\section{Introduction}
\label{ch:intro}

The existence of minimal hypersurfaces is a central problem in differential geometry and the calculus of variations. The Almgren--Pitts min-max theory~\cite{Alm65, Pit81}, together with the regularity results of Schoen--Simon~\cite{SS81}, establishes that every closed Riemannian manifold $M^{n+1}$ with $2 \leq n \leq 6$ contains a smooth, closed, embedded minimal hypersurface. This foundational result has since inspired far-reaching developments, including the resolution of the Willmore conjecture by Marques--Neves~\cite{MN14}, Yau's conjecture on infinitely many minimal hypersurfaces by Marques--Neves~\cite{MN17} (for positive Ricci curvature) and Song~\cite{Song18} (in full generality), and the generic equidistribution results of Irie--Marques--Neves~\cite{IMN17} and Marques--Neves--Song~\cite{MNS}.

Two natural directions of generalization arise. The first concerns manifolds \emph{with boundary}: given a compact Riemannian manifold $M$ with non-empty boundary $\partial M$, one seeks \emph{free boundary minimal hypersurfaces} --- hypersurfaces that are minimal in the interior and meet $\partial M$ orthogonally along their boundary. Li--Zhou~\cite{LZ17} developed a min-max theory for free boundary minimal hypersurfaces and proved that every smooth compact Riemannian manifold $M^{n+1}$ ($2 \leq n \leq 6$) with $\partial M \neq \varnothing$ contains such a hypersurface, without any assumption on the geometry or topology of $M$ or $\partial M$.

The second direction concerns \emph{non-compact} manifolds. For complete non-compact manifolds without boundary, Chambers--Liokumovich~\cite{CL20} introduced a localized min-max framework based on \emph{good sets}: bounded open sets $U \subset M$ whose interior boundary $\partial_I U$ is small relative to the width $W(U)$. The intuition is as follows (see Figure~\ref{fig:good-set-intro}): on a manifold of finite volume (or, more generally, of sublinear volume growth), one can find a bounded open set $U$ whose interior boundary $\partial_I U$ is located at a narrow neck, so that $\mathbf{M}(\partial_I U)$ is small compared to the width $W(U)$. This smallness is crucial for the min-max construction: the initial and terminal slices of the sweepout can be chosen to have mass bounded by $\mathbf{M}(\partial_I U)$, so the min-max value --- the width $W(U)$ --- is strictly larger than the mass of the boundary slices, guaranteeing that the min-max procedure produces a non-trivial minimal hypersurface. Using good sets, Chambers--Liokumovich proved that every complete non-compact manifold of finite volume (more generally, of sublinear volume growth) contains a closed embedded minimal hypersurface.

This thesis combines both directions: we develop a min-max theory for free boundary minimal hypersurfaces on \emph{complete, non-compact} Riemannian manifolds \emph{with boundary}. Our main result is the following.

\begin{figure}[ht]
\centering
\begin{tikzpicture}[scale=0.8]
  %% (b) 3D symmetric bottle — dM transverse to neck
  \begin{scope}[xshift=3.5cm]
    % Fill body only (no tube fill above ∂_I U)
    \shade[left color=fillB!60, right color=fillB!15]
      (-1.5,0)
      .. controls (-1.5,0.3) and (-2,0.7) .. (-2,1.5)
      .. controls (-2,2.8) and (-0.5,3.2) .. (-0.5,3.5)
      arc (180:360:0.5 and 0.13)
      .. controls (0.5,3.2) and (2,2.8) .. (2,1.5)
      .. controls (2,0.7) and (1.5,0.3) .. (1.5,0)
      arc (0:-180:1.5 and 0.4);

    % Left profile
    \draw[thick] (-1.5,0)
      .. controls (-1.5,0.3) and (-2,0.7) .. (-2,1.5)
      .. controls (-2,2.8) and (-0.5,3.2) .. (-0.5,3.5)
      -- (-0.3,5.5);
    % Right profile
    \draw[thick] (1.5,0)
      .. controls (1.5,0.3) and (2,0.7) .. (2,1.5)
      .. controls (2,2.8) and (0.5,3.2) .. (0.5,3.5)
      -- (0.3,5.5);

    % Base ellipse — front half (solid, ∂M)
    \draw[very thick] (1.5,0) arc (0:-180:1.5 and 0.4);
    % Base ellipse — back half (dashed, light blue = ∂_F U)
    \draw[thick, dashed, bindB!40] (1.5,0) arc (0:180:1.5 and 0.4);
    \node[below] at (0,-0.55) {$\partial M = \textcolor{bindB!40}{\partial_F U}$};

    % Neck ellipse (∂_I U) — front half
    \draw[thick, bindB] (0.5,3.5) arc (0:-180:0.5 and 0.13);
    % Neck ellipse — back half (dashed)
    \draw[thick, bindB, dashed] (0.5,3.5) arc (0:180:0.5 and 0.13);
    \node[bindB, right] at (0.65,3.5) {$\partial_I U$};

    % Label U
    \node[bindB] at (0,1.3) {\large $U$};

    % Dots for non-compactness
    \node at (-0.3,5.8) {$\vdots$};
    \node at (0.3,5.8) {$\vdots$};
    \node[above, font=\small] at (0,6.0) {$\to \infty$};

    % Narrow neck annotation
    \draw[<-, >=stealth, thin] (0.5,3.9) -- (1.8,4.5)
      node[right, font=\small] {narrow neck};

    % Closed minimal hypersurface (red ellipse at widest point)
    \draw[thick, bindA] (2,1.5) arc (0:-180:2 and 0.5);
    \draw[thick, bindA, dashed] (2,1.5) arc (0:180:2 and 0.5);
    \node[bindA, right, font=\small] at (2.1,1.5) {$\Sigma$};

    % ∂_F U overlay on front arc (topmost layer, light blue)
    \draw[very thick, bindB!40] (1.5,0) arc (0:-180:1.5 and 0.4);

    % Subcaption
    \node[below, font=\small] at (0,-1.3) {(b)};
  \end{scope}

  %% (a) 3D half-bottle — dM = cut face, parallel to neck
  \begin{scope}[xshift=-3.5cm]
    % Fill body (full cross-section: right profile → neck back arc → back cut face)
    \shade[left color=fillB!60, right color=fillB!15]
      (0,-0.65)
      .. controls (1.5,-0.65) and (2,0.3) .. (2,1.5)
      .. controls (2,2.8) and (0.5,3.2) .. (0.5,3.5)
      arc (0:100:0.5 and 0.13)
      .. controls (-0.087,3.0) and (-0.35,2.2) .. (-0.35,1.523)
      arc (-180:-90:0.35 and 2.173) -- cycle;

    % Right profile (outer surface, from bottom vertex up)
    \draw[thick] (0,-0.65)
      .. controls (1.5,-0.65) and (2,0.3) .. (2,1.5)
      .. controls (2,2.8) and (0.5,3.2) .. (0.5,3.5)
      -- (0.35,5.5);


    % Tube of cut face — front line (from front edge top to infinity)
    \draw[thick] (0.087,3.372) -- (0.06,5.5);
    % Tube of cut face — back line (from back edge top to infinity)
    \draw[thick] (-0.087,3.628) -- (-0.06,5.5);

    % Cut face — front side: bottom arc + top bezier
    \draw[thick] (0,-0.65) arc (-90:0:0.35 and 2.043)
      .. controls (0.35,2.5) and (0.087,3.1) .. (0.087,3.372);

    % Cut face — back side: bottom arc + top bezier
    \draw[thick] (0,-0.65) arc (-90:-180:0.35 and 2.173)
      .. controls (-0.35,2.2) and (-0.087,3.0) .. (-0.087,3.628);

    % Neck (∂_I U) — front arc (solid, slightly shortened)
    \draw[thick, bindB] (0.087,3.372) arc (-80:0:0.5 and 0.13);
    % Neck — back arc (dashed, slightly extended)
    \draw[thick, bindB, dashed] (0.5,3.5) arc (0:100:0.5 and 0.13);
    \node[bindB, right] at (0.65,3.5) {$\partial_I U$};

    % Label U
    \node[bindB] at (1.2,2.3) {\large $U$};

    % Labels for ∂M and ∂_F U
    \node[left, font=\small] at (-0.15,4.8) {$\partial M$};
    \node[bindB!40, below left, font=\small] at (-0.3,-0.65) {$\partial_F U$};

    % Dots for non-compactness (one above each vertical line)
    \node at (-0.06,5.8) {$\vdots$};
    \node at (0.06,5.8) {$\vdots$};
    \node at (0.35,5.8) {$\vdots$};
    \node[above, font=\small] at (0.15,6.0) {$\to \infty$};

    % Narrow neck annotation
    \draw[<-, >=stealth, thin] (-0.5,3.9) -- (-1.8,4.5)
      node[left, font=\small] {narrow neck};

    % Free boundary minimal hypersurface (red ellipse, same size as figure (b))
    % Front solid arc: stop at cut face front surface (θ=-80°, x≈0.35)
    \draw[thick, bindA] (2,1.5) arc (0:-80:2 and 0.5);
    % Back dashed arc: stop at cut face back surface (θ=100°, x≈-0.35)
    \draw[thick, bindA, dashed] (2,1.5) arc (0:100:2 and 0.5);
    \node[bindA, right, font=\small] at (2.1,1.5) {$\Sigma$};

    % ∂_F U overlay on front cut face (topmost layer, light blue)
    \draw[very thick, bindB!40] (0,-0.65) arc (-90:0:0.35 and 2.043)
      .. controls (0.35,2.5) and (0.087,3.1) .. (0.087,3.372);
    % ∂_F U overlay on back cut face (topmost layer, light blue)
    \draw[very thick, bindB!40] (0,-0.65) arc (-90:-180:0.35 and 2.173)
      .. controls (-0.35,2.2) and (-0.087,3.0) .. (-0.087,3.628);

    % Intersection dots with ∂_F U (topmost layer)
    \fill[bindA] ({2*cos(-80)},{1.5+0.5*sin(-80)}) circle (2.5pt);
    \fill[bindA] ({2*cos(100)},{1.5+0.5*sin(100)}) circle (2.5pt);
    % ∂Σ label with arrows to both dots
    \node[bindA, left, font=\small] (dSigma) at (-1.2,1.5) {$\partial\Sigma$};
    \draw[->, >=stealth, bindA, thin, shorten >=4pt] (dSigma) -- ({2*cos(-80)},{1.5+0.5*sin(-80)});
    \draw[->, >=stealth, bindA, thin, shorten >=4pt] (dSigma) -- ({2*cos(100)},{1.5+0.5*sin(100)});

    % Subcaption
    \node[below, font=\small] at (0.5,-1.0) {(a)};
  \end{scope}
\end{tikzpicture}
\caption{A good set $U$ (shaded) inside a non-compact manifold $M$ with boundary $\partial M$. The \emph{interior boundary} $\partial_I U$ sits at a narrow neck, so $\mathbf{M}(\partial_I U)$ is small relative to $W_\partial(U)$; the \emph{free boundary} $\partial_F U = \overline{U}\cap\partial M$ is shown in light blue. The min-max hypersurface $\Sigma$ (red) produced by Theorem~\ref{thm:main} may or may not have boundary, depending on the geometry of $\partial M$: (a)~$\partial M$ runs parallel to the neck, and $\Sigma$ is a \emph{free boundary} minimal hypersurface with $\partial\Sigma\neq\varnothing$ meeting $\partial M$ orthogonally (red dots); (b)~$\partial M$ is transverse to the neck, and $\Sigma$ is a \emph{closed} minimal hypersurface ($\partial\Sigma=\varnothing$) contained in the interior of~$M$.}
\label{fig:good-set-intro}
\end{figure}

\begin{theorem}[Main Theorem]\label{thm:main}
Let $(M^{n+1}, g)$ be a complete Riemannian manifold with non-empty smooth boundary $\partial M$, $2 \leq n \leq 6$.  Let $U \subset M$ be a \emph{good set} (Definition~\ref{def:good-set}), i.e., a bounded open set of finite perimeter satisfying $\mathbf{M}(\partial_I U) \leq \frac{1}{10}\, W_\partial(U)$, with $\overline{U} \cap \partial M \neq \varnothing$.  Then there exists a smooth, \underline{almost properly embedded} minimal hypersurface $(\Sigma,\partial\Sigma)\subset(M,\partial M)$ with (possibly empty) boundary $\partial\Sigma\subset\partial M$ such that
\begin{enumerate}[label=\textup{(\roman*)}]
    \item $\Sigma$ is minimal in the interior and meets $\partial M$ orthogonally along $\partial\Sigma$;
    \item \textup{(Area bound)} $\mathcal{H}^n(\Sigma)\le W_\partial(U)+\mathbf{M}(\partial_IU)$;
    \item \textup{(Localization)} $\Sigma\cap\overline{U}\neq\varnothing$.
\end{enumerate}
\end{theorem}

We now explain the conclusions of the theorem. First, $\Sigma$ is asserted to be \emph{almost properly embedded} in $M$, a notion introduced by Li--Zhou~\cite{LZ17}. This means that $\Sigma \subset M$ and $\partial\Sigma \subset \partial M$, but $\Sigma$ may additionally touch $\partial M$ tangentially at interior points of $\Sigma$. Such points are called \emph{false boundary points}: at a false boundary point $p \in (\Sigma \setminus \partial\Sigma) \cap \partial M$, the tangent plane of $\Sigma$ coincides with that of $\partial M$, and $\Sigma$ touches $\partial M$ from the interior of $M$ (see \cite[Figure~1]{LZ17} for an illustration). The stronger condition that $\Sigma \cap \partial M = \partial\Sigma$ (no false boundary points) is called \emph{proper embedding}. This subtlety does not arise in the boundaryless case of Chambers--Liokumovich~\cite{CL20}, where the min-max hypersurface is embedded in the usual sense. In the presence of $\partial M$, however, the regularity theory (\cite[Theorem~6.2]{LZ17}; see Section~\ref{ch:convergence-of-min-max}) relies on a compactness theorem for stable free boundary minimal hypersurfaces (\cite[Theorem~2.13]{LZ17}), which only preserves almost proper embedding: a sequence of properly embedded hypersurfaces may converge to a limit that is merely almost properly embedded, acquiring false boundary points in the limit (see \cite[Figure~2]{LZ17}). We note that this issue is confined to the regularity theory: the sweepout construction, nested sweepout machinery, and no-escape-to-infinity argument (Sections~\ref{ch:morse-foliation}--\ref{ch:no-escape-to-infinity}) operate entirely at the level of currents and do not involve smooth hypersurfaces, so false boundary points play no role there.

Turning to the three numbered conclusions: (i)~is the defining property of a free boundary minimal hypersurface, as introduced in~\cite{LZ17}. In~(ii), $W_\partial(U)$ denotes the \emph{relative width} of $U$ (Definition~\ref{def:relative-sweepout}), which is the width computed from sweepouts whose slices are restricted to $\overline{U}$; it measures the intrinsic ``waist'' of $U$ (see Remark~\ref{rem:why-relative-width} for why this is the appropriate notion of width).

Conclusions~(ii) and~(iii) together address issues specific to the non-compact setting. The localization condition $\Sigma \cap \overline{U} \neq \varnothing$ ensures that the limit hypersurface does not escape to infinity: it genuinely passes through the good set $U$. However, localization alone does not prevent the hypersurface from extending far beyond $U$ with most of its area at infinity; the area bound provides the complementary control. Combined with the good set condition $\mathbf{M}(\partial_I U) \leq \frac{1}{10}\, W_\partial(U)$ (Definition~\ref{def:good-set}), it yields $\mathcal{H}^n(\Sigma) \leq \frac{11}{10}\, W_\partial(U)$, so the area of $\Sigma$ is controlled at the scale of the width of $U$, giving a geometrically meaningful minimal hypersurface that lives at the scale of the good set (see Figure~\ref{fig:good-set-intro}). In the compact setting both conclusions are automatic, but on non-compact manifolds they require the good set machinery of Chambers--Liokumovich~\cite{CL20} together with the no-escape-to-infinity argument developed in Section~\ref{ch:no-escape-to-infinity}.

\begin{remark}[Why the relative width]\label{rem:why-relative-width}
Several notions of width appear in the min-max framework: the (unrestricted) width $W(U)$, the \emph{relative width} $W_\partial(U)$ (Definition~\ref{def:relative-sweepout}), and the \emph{good width} $W_g(U)$ (Definition~\ref{def:good-sweepout}). In the boundaryless case, Chambers--Liokumovich~\cite{CL20} establish
\[
W_\partial(U) \;\leq\; W(U) = W_g(U) \;\leq\; W_\partial(U) + \mathbf{M}(\partial U),
\]
where $\partial U$ is the full boundary of $U$ inside $M$. In our setting, where $M$ has boundary $\partial M$, the boundary of $U$ splits as $\partial U = \partial_I U \cup \partial_F U$, and only the interior boundary $\partial_I U$ contributes to the gap; we prove (Proposition~\ref{prop:width-inequalities}) that the same inequalities hold with $\mathbf{M}(\partial U)$ replaced by $\mathbf{M}(\partial_I U)$:
\[
W_\partial(U) \;\leq\; W(U) = W_g(U) \;\leq\; W_\partial(U) + \mathbf{M}(\partial_I U).
\]
Here $W(U) = W_g(U)$ is a non-trivial theorem (Proposition~\ref{prop:width-inequalities}(iii)). The gap $W(U) - W_\partial(U) \leq \mathbf{M}(\partial_I U)$ arises because slices of an unrestricted sweepout may spill over $\partial_I U$, contributing extra area up to $\mathbf{M}(\partial_I U)$ that is irrelevant to the geometry of $U$ itself. The good width $W_g(U)$ plays a key role in the \emph{proof} --- the pulled-tight process requires good sweepouts whose initial and terminal slices have small mass --- but in the \emph{statement} of the theorem, the relative width $W_\partial(U)$ gives the sharpest area bound.
\end{remark}

\medskip
\noindent\textbf{Organization of the thesis.}
Section~\ref{ch:preliminaries} collects the analytical foundations: integer rectifiable currents, relative cycles, varifolds, and the sweepout framework needed for the min-max construction on manifolds with boundary.

The core of the proof occupies Sections~\ref{ch:morse-foliation}--\ref{ch:convergence-of-min-max}. Section~\ref{ch:morse-foliation} develops Morse foliations with controlled fiber area, the building blocks for converting arbitrary sweepouts into nested ones. Section~\ref{ch:splitting-and-extension} establishes the splitting and extension lemmas that allow local nested families to be glued together with controlled mass increase. These tools are combined in Section~\ref{ch:nested-sweepout}, which constructs nested sweepouts from arbitrary sweepouts --- the key technical input for the no-escape argument. Section~\ref{ch:no-escape-to-infinity} proves that any good sweepout must carry definite mass inside the good set $U$, preventing the min-max sequence from escaping to infinity. Finally, Section~\ref{ch:convergence-of-min-max} completes the proof: a pull-tight procedure adapted to the non-compact setting via exhaustion produces a stationary varifold, which is shown to be almost minimizing in small annuli with free boundary; the regularity theory of Li--Zhou~\cite{LZ17} then yields the smooth, almost properly embedded free boundary minimal hypersurface of Theorem~\ref{thm:main}.